\section{OpenMP--MPI Implementation}

The CPU-parallel backend follows the natural coarse-grained decomposition of ACO.  OpenMP threads build independent ant solutions inside each rank, while MPI combines information across ranks.  This mapping keeps the sequential decision chain of an individual ant local and exposes parallelism across ants.

\subsection{Shared-memory organization}

The current implementation creates one persistent OpenMP region around the epoch loop.  Thread-private route-construction workspaces avoid sharing the visited set and random-number state.  Candidate indices and heuristic scores are stored contiguously, while the distance and pheromone matrices use single-precision aligned storage within the parallel solver.  Per-epoch work consists of four main phases: candidate-score refresh, independent ant construction, selection of the best solution, and pheromone evaporation/deposit.

The ant loop uses a runtime-selectable OpenMP schedule.  Pheromone evaporation is partitioned statically over the dense matrix; deposits use atomic additions because different routes can touch the same edge.  This removes a serial deposit phase at the price of possible atomic contention.  The fixed-size strong-scaling campaign in Section~\ref{sec:performance} is designed to measure the net effect rather than attributing time to any one phase without profiling.

\subsection{Distributed-memory synchronization}

Each MPI rank owns a replica of the pheromone matrix and receives a disjoint portion of the total ant population.  Modified edges are represented as sparse \emph{delta} records containing a matrix index and an increment.  The implementation starts a nonblocking \texttt{MPI\_Iallgatherv}; deltas from the previous epoch are applied before the next local update.  This avoids transmitting the complete dense matrix when only route edges have changed.

The sparse exchange does not remove all distributed overhead.  Every rank still initializes and evaporates a dense matrix, performs collective coordination, and stores its local replica.  Consequently, increasing the number of ranks may expose replicated work even when the communicated payload is sparse.  The observed multi-node trend is discussed conservatively in Section~\ref{sec:performance}; no network or phase-level profile is available in the current result set.

\subsection{Alternatives that were rejected}

The first CPU-parallel version exposed four concrete implementation problems:
\begin{itemize}[leftmargin=*,nosep]
    \item \textbf{False sharing:} adjacent, unaligned per-thread telemetry structures generated unnecessary cache-coherence traffic.
    \item \textbf{Repeated fork--join:} opening and closing an OpenMP region at every epoch added recurring thread-management overhead.
    \item \textbf{Serial deposits:} pheromone deposits executed inside a single-thread section and serialized part of every epoch.
    \item \textbf{Fine-grained MPI synchronization:} one nonblocking reduction per matrix row produced a large number of small collectives.
\end{itemize}

These problems motivated the organization retained in the submitted version: one persistent OpenMP region, aligned thread-private workspaces, atomic parallel deposits, and sparse delta exchange through \texttt{MPI\_Iallgatherv}.  The changes remove the original serial and per-row synchronization structures, but do not make coordination free.  Atomic contention, phase barriers, replicated dense-matrix work, and collective synchronization remain relevant when the ant population does not provide enough independent work per worker.

Further alternatives exposed the following limitations:
\begin{itemize}[leftmargin=*,nosep]
    \item \textbf{Scheduling overhead and imbalance:} fine-grained dynamic scheduling repeatedly accesses the OpenMP work queue, whereas larger guided chunks can leave a long tail because ants construct routes of different lengths.
    \item \textbf{Insufficient task granularity:} a fixed population may expose fewer independent ants than threads.  Once every ant has been assigned, extra workers can only assist the dense matrix phases; this limitation remains visible in the strong-scaling campaign.
    \item \textbf{Dense-matrix pressure:} double precision enlarged the working set and memory traffic.  Single precision reduces this cost but preserves quadratic storage.
    \item \textbf{Irregular fallback:} if the candidate list has no feasible unvisited customer, selection falls back to an $O(n)$ scan.  Two-level candidate lists, hierarchical visited masks, and cache-adaptive candidate widths reduce exhausted-block work, but cannot remove the global fallback without changing the search rule.
    \item \textbf{Poor prefetch and SIMD fit:} the visited-dependent access pattern frustrates manual prefetching, while direct SIMD requires masked filtering, reduction, and probabilistic selection and can evaluate entries that are immediately discarded.
    \item \textbf{Fine-grained intra-ant coordination:} per-step barriers repeatedly synchronized workers; atomic spinning replaced barrier calls with cache-line bouncing; and condition-variable wakeups introduced operating-system latency comparable to the small scan being parallelized.
\end{itemize}

The final one-ant-per-worker mapping therefore synchronizes only at phase boundaries.  Hierarchical masks and adaptive candidate data are retained, while the rejected implementation history is reported qualitatively and no cross-version speedup is reused.

\subsection{Correctness boundary}

The current validation reconstructs every saved route and checks customer coverage and capacity.  All available OpenMP--MPI routes are feasible.  For several multi-rank runs, however, the scalar \texttt{best\_cost} does not match the cost of the route printed by rank zero.  This does not change the fixed-epoch wall time, which is retained for scaling, but it prevents using those cost values as multi-rank quality evidence.  Resolving the ownership and transfer of the globally best route is therefore a necessary correctness improvement beyond the timing study.
